\section{An explicit uniform bound on the principal Eisenstein component}
\label{sec:explicit-family}

The next statements supply a quantitative version of T5 on a precisely
normalized component. Let $(A,\theta,\eta)$ be a principally polarized
complex abelian fourfold with an integral action
$\eta:\mathcal O_K\to\End(A)$, $K=\Q(\omega)$, of signature $(2,2)$,
and assume the Rosati involution induces complex conjugation on $K$.
The Weil plane $W$ is primitive of type $(2,2)$. Indeed, cup product with
$\theta$ introduces both $K$ characters, so vanishes on the two extreme
exterior powers. The Hodge--Riemann form on $W$ is therefore the positive
intersection form. These are cohomological statements about the projector
motive $(A,e_W)$, not assertions that an arbitrary effective representative
deforms from a single special fiber.

\begin{proposition}[R10: bounded detection gives an integral Weil basis]
Suppose an effective surface $Z$ has $\deg_\theta Z\leq D_0$ and
$e_W[Z]\ne0$. Then $u=54e_W[Z]$ and $Ru$, where $R=\eta(\omega)^*$,
are an integral basis over $\Q$ of $W$, each with a signed effective
representative having both numerator degrees at most $36D_0$. If
$d=(u,u)$, their Gram matrix is
\[
 d\begin{pmatrix}1&-1/2\\-1/2&1\end{pmatrix},
 \qquad 0<d\leq729D_0^2,
 \qquad d\in2\Z.
\]
\end{proposition}
\begin{proof}
The projector calculation P06 applies to every such fourfold. Put
$S=\eta(\sqrt{-3})^*$. Expansion gives the integral representative
\begin{align*}
 u={}&[18Z+2\eta(\sqrt{-3})^*Z]\\
 &-[9\eta(\omega)^*Z+9\eta(\omega^2)^*Z
       +\eta(\sqrt{-3}\omega)^*Z+
        \eta(\sqrt{-3}\omega^2)^*Z].
\end{align*}
All the maps are isogenies and hence flat, so these pullbacks are effective.
The unit maps preserve $\theta$. The map $\eta(\sqrt{-3})$ has
polarization multiplier $3$ and degree $3^4$, so its codimension-two
pullback multiplies $\theta$-degree by $3^4/3^2=9$. Both displayed sides
therefore have degree at most $36D_0$, and the same is true after $R$.

The Rosati involution is the adjoint for the polarized Hodge metric.
Consequently the commuting $K$ actions on the exterior algebra are normal,
their character spaces are orthogonal, and $e_W$ is an orthogonal
projector. Use the invariant flat K\"ahler metric representing $\theta$;
its total volume is $\theta^4/4!=1$. Harmonic forms are translation
invariant, so a unit $L^2$ harmonic form has pointwise Euclidean norm one
and comass at most one. If $\alpha$ is the harmonic representative of
$[Z]$, Poincar\'e duality gives
\[
 |\langle\alpha,\beta\rangle|
   =\left|\int_Z *\beta\right|
   \leq\operatorname{vol}(Z)=\deg_\theta Z/2
\]
for every unit harmonic four-form $\beta$. Singularities are harmless:
integrate over the smooth locus with the cycle multiplicities. Hence
$\|\alpha\|\leq D_0/2$ and
$d=\|54e_W\alpha\|^2\leq729D_0^2$.

On $W$, $R$ is an isometry with $1+R+R^2=0$. If $e=(u,Ru)$, invariance
and symmetry give $(u,R^2u)=e$, so $d+2e=0$. Since $u,Ru$ are integral,
$e$ is integral and $d$ is even. The positive Gram determinant is
$3d^2/4$, proving independence and the asserted matrix.
\end{proof}

\begin{theorem}[R11: explicit extraction cost]
Let $t\ne0$ be a rational vector of $W$, choose $\delta\in\Z_{>0}$
such that $\delta t$ is integral, and put $T=t^2>0$. Under the hypothesis
of R10, the integer
\[
 C_4(\delta,T;D_0)=531441\,\delta D_0^4
              \bigl(1+4\lceil\sqrt T\rceil\bigr)
 \tag{R11.1}
\]
bounds the multiplier and the two effective numerator degrees of an
algebraic representative of $t$.
\end{theorem}
\begin{proof}
Set $M=729D_0^2$, $p=(\delta t,u)$, $q=(\delta t,Ru)$ and
$\Delta=3d^2/4$. These last three numbers are integers and $\Delta>0$.
Inverting the Gram matrix gives
\[
 \Delta\delta t=a u+bRu,\qquad
 a=dp+\frac d2q,\quad b=\frac d2p+dq.
\]
The coefficients are integers. Positive-definite Cauchy--Schwarz gives
$|p|,|q|\leq\delta\sqrt{Td}$ and
$|a|+|b|\leq3\delta\sqrt T\,d^{3/2}$.
Take signed sums of the two representatives from R10. The resulting
multiplier and degrees satisfy
\[
 N=\delta\Delta\leq\delta M^2,\qquad
 \deg Z^\pm\leq108\delta D_0\sqrt T\,M^{3/2}.
\]
Now $729^2=531441$ and $108\cdot729^{3/2}=4\cdot531441$.
The integer in (R11.1) exceeds both bounds. This argument explicitly
identifies the division: the inverse integral intersection matrix on
the two Weil generators is cleared by $\Delta$, followed by the
denominator $\delta$ of the desired class.
\end{proof}

\begin{theorem}[R12: an explicit detecting-degree bound]
On the principal integral Eisenstein component of signature $(2,2)$,
put
\[
 D_0=132269760.
\]
For every fixed flat rational Weil class $t$, the bound (R11.1) holds
throughout that component after trivializing its finite Weil monodromy.
It therefore holds at every good geometric reduction of every arithmetic
member, with the same extended polarization and class. This is an upper
bound; it is not the least complexity at a general member.
\end{theorem}
\begin{proof}
We first construct a bounded detector on a dense set. Take a connected
unramified cyclic triple cover $\pi:C\to X$ with $g(X)=3$ and $g(C)=7$.
Write $J=\operatorname{Jac}(C)$ with its principal $\Theta$, and let
$B$ be the connected Prym, of dimension four. Put $L=\Theta|_B$.
Schoen's actual Prym-dominance calculation, not the overbroad original
formulation of his theorem, is the input: \cite[\S3, pp.~26--30]{Schoen1988}.
The polarization qualification is explained in
\cite[p.~330]{Schoen1998}.

Here is the needed polarization calculation. Let $Q=\im(\pi^*)\subset J$.
The map $\operatorname{Jac}(X)\to Q$ has kernel cyclic of order three,
and pulls $\Theta|_Q$ back to $3\Theta_X$. Thus its polarization has
degree $3^6/3^2=81$. Complementarity in the principally polarized $J$
identifies the two restriction kernels with $B\cap Q$, so $|K(L)|=81$.
The norm operator is zero on $B$ and is multiplication by three on $Q$;
therefore $K(L)$ is killed by three. It follows that
$K(L)\simeq(\Z/3)^4$, that $L$ has type $(1,1,3,3)$, and that $L^4=216$.
These complement and descent identities are the usual polarization
identities for dual isogenies; see \cite[\S11, pp.~17--18; Theorem13.3,
p.~25; Proposition16.8 and Remark16.9, p.~33]{MilneAV}. A direct kernel
and connected-complement justification, valid in both genera used here,
is given in the proof of R20 below.

Schoen's $r=0$ construction in \cite[Theorem~2.0 and pp.~12--13,
paragraph immediately before (2.3)]{Schoen1988} represents the primitive
character subspace by combinations of components lying over
$|K_X|=\Pp^2\subset\operatorname{Sym}^4X$.
The map
\[
 f:C^4\longrightarrow B,\qquad
 f=P\circ\operatorname{Abel},\quad
 P=3\id-(1+\sigma+\sigma^2),
\]
is dominant and generically finite; the pullback of the Weil plane is
exactly that character subspace. Hence $f_*f^*$ is multiplication by
the nonzero degree of $f$, and some reduced source component has image
of dimension two with nonzero Weil projection after pushforward.

Its Abel image lies in $W_4(C)\cap T_0$, where $T_0$ is a translate
of $B$ in the norm fiber over $K_X$. Every effective divisor representing
a point of that intersection maps to the canonical linear system.
The map $\operatorname{Sym}^4C\to\operatorname{Sym}^4X$ is finite,
so the inverse image of that linear system has dimension two.
Consequently the selected two-dimensional Abel image is an irreducible
component $Z$ of $W_4(C)\cap T_0$. On $T_0$, $P$ is $[3]$ followed by
translation. Both induce invertible rational maps on the relevant
cohomology; thus the reduced component $Z$ itself has nonzero Weil
projection after translation to $B$. No ordered-product multiplicity
or degree of $[3]$ is needed after selecting that reduced component.

Embed $J$ by $3\Theta$, using \cite[Theorem7.3, p.~11]{MilneAV}.
The Poincar\'e formula $[W_4(C)]=\Theta^3/3!$
\cite[p.~149, formulas (6)--(7)]{GriffithsPoincare1982} gives
\[
 \deg_{3\Theta}W_4(C)=3^4\frac{7!}{3!}=68040,
 \qquad \deg_{3\Theta}T_0=3^4\,216=17496.
\]
The general projective B\'ezout inequality, valid for the reduced
components of an improper intersection, bounds the degree of $Z$ by
the product; see \cite[Example~8.4.6, p.~146]{Fulton}.
Since $Z$ is a surface,
\[
 \deg_L Z\leq\frac{68040\cdot17496}{3^2}=132269760.
\]

The deck transformation is the identity on $Q$, hence on $K(L)$.
Choose a maximal isotropic subgroup of $K(L)$ of order nine. It is
$\mathcal O_K$-stable. Its quotient $q:B\to A_0$ carries a principal
polarization $\theta_0$ with $q^*\theta_0=L$. Then $q_*Z$ is effective,
has degree at most $D_0$, and has nonzero Weil projection, by projection
formula and rational invertibility of the isogeny. This quotient
construction is a finite isogeny correspondence on moduli, so its image
still has dimension four.

We identify its principal component explicitly. For an integral
alternating polarization $E$, the associated Hermitian form is
\[
 h(x,y)=E(\omega x,y)-\bar\omega E(x,y),\qquad
 E(x,y)=\Tr_{K/\Q}\bigl(h(x,y)/\sqrt{-3}\bigr).
\]
Principalness makes this an integral unimodular Eisenstein Hermitian
lattice of signature $(2,2)$. Its underlying rational quadratic form
has dimension eight and signature $(4,4)$, so is isotropic by Meyer's
theorem \cite[Corollary~3.3.4]{PetersSterk}. Because $\mathcal O_K$ is
a PID, obtain a primitive integral isotropic vector $e$ by clearing
denominators and content. Unimodularity supplies $f$ with $h(e,f)=1$.
Replacing $f$ by $f+ae$ changes its norm by $\Tr(a)$; since
$\Tr(\mathcal O_K)=\Z$, make $f$ isotropic. This splits a unimodular
hyperbolic plane integrally. The remaining rank-two lattice has
signature $(1,1)$ and determinant $-1$. Its rational diagonal form
is $\langle a,-1/a\rangle$, with isotropic vector $(1,a)$; the same
primitive-vector argument splits it. Hence the lattice is $U\oplus U$,
also the lattice of the product $E^4$ with mixed action. The associated
period domain is connected, so after finite level covers the dense
Prym-quotient image belongs to the component in the theorem.

By R13a, a cover of degree at most six trivializes the Weil monodromy.
Fix $t$ on that cover; its $\delta$ and $T$ are constant. Apply R10--R11
on the dense Prym-quotient image. The bounded-class locus is closed by
R13b, so contains the whole component. Finally B05 descends the complex
cycles of an arithmetic fiber without increasing complexity, and B02
specializes them to every good geometric reduction. The upper bound
uses actual cycles and consequently holds simultaneously in all
compatible realizations.
\end{proof}

The constant does not cover uncharged changes of polarization or of
the integral order. On a fixed polarized Hecke orbit inside this
component it does supply the bounded-division gate of Route~2. The
cycles whose classes were isolated are surfaces from the Prym
construction; the operation that divides their classes is R11's
integral Gram-matrix calculation. The exact value $16$ at the product
point is much smaller than this uniform certificate.
